\documentclass{beamer}

\usepackage{amsmath,amsfonts,amssymb}
\usepackage{graphicx}
\usepackage{physics}
\usepackage{tikz}

\title{Reinforcement Learning for Quantum Control}
\subtitle{A Physics-Oriented Introduction with Applications to Quantum Sensors}
\author{}
\date{}

\begin{document}

%------------------------------------------------
\frame{\titlepage}

%------------------------------------------------
\begin{frame}{Motivation}

Many problems in modern physics can be viewed as optimization problems:

\begin{itemize}
\item finding ground states (variational methods)
\item optimizing control pulses
\item maximizing measurement precision
\end{itemize}

However:

\begin{itemize}
\item objective functions may be unknown
\item gradients may be unavailable or noisy
\item decisions may be sequential
\end{itemize}

Reinforcement learning (RL) provides a framework for
\textbf{learning optimal strategies through interaction}.

\end{frame}

%------------------------------------------------
\begin{frame}{What is Reinforcement Learning?}

Reinforcement learning studies how an agent learns to act
in an environment in order to maximize cumulative reward.

Key ingredients:

\begin{itemize}
\item \textbf{State} $s_t$: description of the system
\item \textbf{Action} $a_t$: decision made by the agent
\item \textbf{Reward} $r_t$: feedback signal
\end{itemize}

The agent interacts with the system over time:

\[
(s_0,a_0,r_0,s_1,a_1,r_1,\dots)
\]

\end{frame}

%------------------------------------------------
\begin{frame}{Sequential Decision Making}

Unlike standard optimization, RL is inherently sequential.

Decisions influence future states:

\[
s_t \rightarrow a_t \rightarrow s_{t+1}
\]

Thus we must consider long-term consequences.

Goal:

\[
\max \mathbb{E}\left[\sum_{t=0}^\infty \gamma^t r_t \right]
\]

This resembles action principles in physics,
where the total trajectory matters.

\end{frame}

%------------------------------------------------
\begin{frame}{Markov Decision Processes}

We assume the system satisfies the Markov property:

\[
P(s_{t+1}|s_t,a_t,\dots)=P(s_{t+1}|s_t,a_t)
\]

An MDP is defined by:

\begin{itemize}
\item states $s$
\item actions $a$
\item transition probabilities $P(s'|s,a)$
\item reward function $R(s,a)$
\end{itemize}

This is analogous to stochastic dynamics in physics.

\end{frame}

%------------------------------------------------
\begin{frame}{Policies}

A policy specifies how actions are chosen:

\[
\pi(a|s)
\]

It can be:

\begin{itemize}
\item deterministic: $a = f(s)$
\item stochastic: probabilities over actions
\end{itemize}

The goal of RL is to find an optimal policy $\pi^*$.

\end{frame}

%------------------------------------------------
\begin{frame}{Value Functions: Physical Interpretation}

Define the value function:

\[
V^\pi(s) = \mathbb{E}_\pi\left[\sum_t \gamma^t r_t \mid s_0=s\right]
\]

Interpretation:

\begin{itemize}
\item expected ``future return''
\item analogous to free energy or action
\end{itemize}

Similarly,

\[
Q^\pi(s,a)
\]

represents the value of taking action $a$ in state $s$.

\end{frame}

%------------------------------------------------
\begin{frame}{Bellman Equation: Self-Consistency}

Key relation:

\[
V^\pi(s) = \sum_a \pi(a|s)\left[
R(s,a) + \gamma \sum_{s'} P(s'|s,a)V^\pi(s')
\right]
\]

Interpretation:

\begin{itemize}
\item immediate reward
\item plus expected future reward
\end{itemize}

This is a recursive equation similar to
self-consistency equations in many-body physics.

\end{frame}

%------------------------------------------------
\begin{frame}{Optimal Control Perspective}

Define optimal value:

\[
V^*(s) = \max_\pi V^\pi(s)
\]

Bellman optimality:

\[
V^*(s) = \max_a \left[
R(s,a) + \gamma \sum_{s'} P(s'|s,a)V^*(s')
\right]
\]

This is the discrete analogue of the
Hamilton--Jacobi--Bellman equation.

\end{frame}

%------------------------------------------------
\begin{frame}{Value Iteration}

We can solve for $V^*$ iteratively:

\[
V_{k+1}(s) = \max_a \left[
R(s,a) + \gamma \sum_{s'} P(s'|s,a)V_k(s')
\right]
\]

This resembles iterative solutions in physics:

\begin{itemize}
\item self-consistent field methods
\item renormalization procedures
\end{itemize}

\end{frame}

%------------------------------------------------
\begin{frame}{Policy Iteration}

Two-step algorithm:

\textbf{1. Policy evaluation}

Compute $V^\pi$

\textbf{2. Policy improvement}

\[
\pi'(s) = \arg\max_a Q^\pi(s,a)
\]

Repeat until convergence.

This is analogous to variational improvement.

\end{frame}

%------------------------------------------------
\begin{frame}{Q-Learning: Model-Free Approach}

In many problems we do not know $P(s'|s,a)$.

Q-learning uses data:

\[
Q(s,a) \leftarrow Q(s,a) + \alpha
\left[r + \gamma \max_a Q(s',a) - Q(s,a)\right]
\]

Interpretation:

\begin{itemize}
\item temporal-difference learning
\item stochastic approximation to Bellman equation
\end{itemize}

\end{frame}

%------------------------------------------------
\begin{frame}{Policy Gradient Methods}

We parameterize the policy:

\[
\pi_\theta(a|s)
\]

Objective:

\[
J(\theta)
\]

Gradient:

\[
\nabla_\theta J =
\mathbb{E}[\nabla_\theta \log \pi_\theta(a|s) Q(s,a)]
\]

This resembles stochastic gradient descent.

\end{frame}

%------------------------------------------------
\begin{frame}{REINFORCE Algorithm}

Monte Carlo estimator:

\[
\nabla_\theta J \approx
\sum_t \nabla_\theta \log \pi(a_t|s_t) G_t
\]

with return:

\[
G_t = \sum_{k=t}^\infty \gamma^{k-t} r_k
\]

Interpretation:

\begin{itemize}
\item increase probability of good trajectories
\item decrease probability of bad ones
\end{itemize}

\end{frame}

%------------------------------------------------
\begin{frame}{Actor--Critic Methods}

Combine:

\begin{itemize}
\item Actor: policy
\item Critic: value function
\end{itemize}

Temporal difference error:

\[
\delta = r + \gamma V(s') - V(s)
\]

Actor update:

\[
\nabla_\theta \log \pi(a|s)\,\delta
\]

This improves stability and efficiency.

\end{frame}

%------------------------------------------------
\begin{frame}{RL as Variational Optimization}

RL can be viewed as optimizing

\[
J(\theta)
\]

but with:

\begin{itemize}
\item unknown objective function
\item stochastic gradients
\item sequential dependence
\end{itemize}

Thus RL generalizes variational methods to
\textbf{adaptive, sequential decision problems}.

\end{frame}

%------------------------------------------------
\begin{frame}{Transition to Quantum Systems}

Many quantum problems naturally fit RL:

\begin{itemize}
\item sequential control
\item measurement feedback
\item stochastic noise
\end{itemize}

We now apply RL to quantum control and sensing.

\end{frame}

%------------------------------------------------
% (Keep quantum control + sensing slides same as before)
%------------------------------------------------

\begin{frame}{Quantum Control}

\[
i \frac{d}{dt}\ket{\psi} = H(t)\ket{\psi}
\]

\[
H(t)=H_0 + \sum_k u_k(t) H_k
\]

Goal: optimize control fields $u_k(t)$.

\end{frame}

%------------------------------------------------
\begin{frame}{RL Mapping}

\begin{itemize}
\item state: quantum state
\item action: control pulse
\item reward: fidelity or Fisher information
\end{itemize}

\end{frame}

%------------------------------------------------
\begin{frame}{Quantum Sensors}

\[
\Delta \theta \ge \frac{1}{\sqrt{\nu F_Q}}
\]

RL can optimize $F_Q$.

\end{frame}

%------------------------------------------------
\begin{frame}{Two-Level Sensor Example}

\[
H = \frac{B}{2}\sigma_z + u_x(t)\sigma_x
\]

Goal: maximize sensitivity to $B$.

\end{frame}

%------------------------------------------------
\begin{frame}{Summary}

RL provides:

\begin{itemize}
\item a general framework for sequential optimization
\item connections to stochastic control and variational methods
\item powerful tools for quantum technologies
\end{itemize}

\end{frame}

\end{document}
